\section{Introduction}
\label{sec:introduction}

High-latency networks are common in satellite communication, non-terrestrial networking, intercontinental transmission, and remote access systems. Compared with conventional terrestrial links, such environments often exhibit long propagation delay, high bandwidth-delay product, variable bandwidth, and non-congestion losses. End-to-end TCP control relies on acknowledgements, round-trip time (RTT) samples, and loss events to adjust its sending behavior. When the RTT is large, however, the feedback loop becomes too slow to track rapid link variations. During the delayed response interval, the sender or proxy may continue injecting traffic according to an outdated state, which can lead to queue buildup, throughput fluctuation, and slow recovery after link switching.

Performance Enhancing Proxies (PEPs) are an engineering approach to mitigate link-related performance degradation by placing local processing functions closer to the bottleneck link. By splitting part of the control loop and introducing local buffering, pacing, and recovery mechanisms, PEPs can reduce the dependence on end-to-end feedback. PEPesc extends this idea by combining a distributed TCP performance enhancing proxy with streaming forward error correction (FEC), which reduces the need for RTT-scale retransmission after losses. Nevertheless, the original PEPesc design is more suitable for static or slowly changing links. Under dynamic bandwidth and loss conditions, the proxy-side estimator may lag behind the reference link state, the control action lacks explicit state semantics, and pacing control is not sufficiently coordinated with FEC redundancy.

This paper focuses on improving the quality of networking service for dynamic high-latency transmission without redesigning a complete end-to-end congestion control protocol. Instead, we enhance the PEPesc proxy-side loop with a perception-decision-execution framework. At the perception layer, we keep the rule-based estimator as a stable baseline and introduce a hybrid residual GRU model to correct bandwidth and loss-rate estimation errors. At the decision layer, we design an ACCEL/HOLD/BRAKE state machine to distinguish link recovery, stable transmission, and congestion suppression. At the execution layer, we apply conservative pacing correction and an FEC reliability floor to avoid unstable adaptation and preserve basic streaming recovery capability.

The main contributions of this paper are summarized as follows:
\begin{itemize}
    \item We propose a GRU-assisted link perception method for PEP-side dynamic link estimation. The method learns residual corrections over a rule-based baseline rather than replacing the baseline estimator entirely.
    \item We design a three-state adaptive control mechanism with explicit ACCEL, HOLD, and BRAKE semantics, enabling interpretable control decisions under dynamic bandwidth, queue, RTT, ACK, and loss conditions.
    \item We implement a conservative execution strategy that combines legacy baseline behavior with adaptive correction and an FEC reliability floor, reducing the risk that adaptive control weakens streaming recovery.
    \item We evaluate the enhanced system using Mininet-based high-latency emulation. The results show improvements in receiver throughput, recovery time, throughput deficit area, p95 queue delay, and long-duration effective running time.
\end{itemize}

The remainder of this paper is organized as follows. Section~\ref{sec:related} reviews related work. Section~\ref{sec:framework} presents the proposed framework. Section~\ref{sec:setup} describes the implementation and experimental setup. Section~\ref{sec:evaluation} reports the evaluation results. Section~\ref{sec:discussion} discusses the findings and limitations, and Section~\ref{sec:conclusion} concludes the paper.
